\documentclass[12pt]{article}

\usepackage[top=1in,bottom=1in,left=1in,right=1in]{geometry}
\usepackage{amsmath,amssymb,mathtools}
\usepackage{xcolor}
\usepackage{graphicx}
\usepackage{hyperref}
\usepackage{enumitem}

\hypersetup{
    pdfauthor={},
    pdftitle={MATH 1075 - Lesson 6.1 Recitation}
}

\begin{document}

\begin{center}
    {\LARGE MATH 1075 - Section 10.6 Worksheet}

    \vspace{2em}

    {\large 11/17/2026}
\end{center}
\section{Simplifying Radicals with Fractions}
Simplify the following radical expressions. (Use $\sqrt[n]{\frac{a}{b}}=\frac{\sqrt[n]{a}}{\sqrt[n]{b}}, b\neq  0 $)

\vspace{1.5em}

1. $\sqrt{\frac{a^{6}}{b^{4}} }$

\vspace{4.5em} 

2. $\sqrt[3]{\frac{125}{8} }$

\vspace{4.5em} 

3. $\sqrt[3]{\frac{-16j^{3}}{k^{3}} }$

\vspace{4.5em} 

4. $\sqrt{\frac{6x^{3}}{24x} }$

\vspace{4.5em}

5. $\frac{\sqrt[4]{3b^{3}}}{\sqrt[4]{48b^{11}}} $ 

\vspace{4.5em}

\section{Radicals in Denominator - One Term} 
First, try to identify what $\sqrt{?}$ is that will make the following equation true, assuming all variables are positive real numbers. This process of finding $\sqrt{?}$ and getting rid of the radicals is called \textbf{rationalizing} (Recall: $\sqrt[n]{a}\cdot\sqrt[n]{b}=\sqrt[n]{ab}$).

\vspace{1.5em} 

1. $\sqrt{a}\cdot \sqrt{?}=\sqrt{a^{2}}=a$ 

\vspace{4.5em} 

2. $\sqrt[3]{b}\cdot \sqrt[3]{?}=\sqrt[3]{b^{3}}=b$ 

\vspace{4.5em} 

3. $\sqrt[4]{2z^{3}}\cdot \sqrt[4]{?}=\sqrt[4]{2^{4}z^{4}}=2z$ 

\vspace{4.5em} 

\noindent Rationalize the Denominator (Remember, if you're multiplying something to the denominator, you also have to do the same for the numerator): 

\vspace{1.5em} 

1. $\frac{2}{\sqrt{a}} $ 

\vspace{4.5em} 

2. $\frac{1}{\sqrt[3]{b}} $ 

\vspace{4.5em} 

3. $\sqrt[4]{\frac{16}{2z^{3}}}$

\vspace{4.5em} 

\section{Radicals in Denominator - Two Terms}
Try to identify what $?$ is that will make the following equation true. (Hint: Recall that we have $(a+b)(a-b)=a^{2}-b^{2}$.)

\vspace{1.5em} 

1. $(4-\sqrt{3})\cdot ?=4^{2}-(\sqrt{3})^{2}=16-3=13$ 

\vspace{4.5em} 

2. $(\sqrt{p}+\sqrt{q})\cdot ?=(\sqrt{p})^{2}-(\sqrt{q})^{2}=p-q$ 

\vspace{4.5em} 

3. $(\sqrt{x}-\sqrt{5})\cdot ?=(\sqrt{x})^{2}-(\sqrt{5})^{2}=x-5$ 

\vspace{4.5em} 
\noindent Rationalize the Denominator (Remember, if you're multiplying something to the denominator, you also have to do the same for the numerator): 

\vspace{1.5em} 

1. $\frac{6}{4-\sqrt{3}}  $ 

\vspace{4.5em} 

2. $\frac{-1}{\sqrt{p}+\sqrt{q}} $ 

\vspace{4.5em} 

3. $\frac{x-5}{\sqrt{x}-\sqrt{5}} $

\vspace{4.5em} 
\end{document}